% Appendix 10: Visual Representation of Sections of Four-Dimensional Space
% Source pages: 284–286

\chapter{Visual Representation of Sections of Four-Dimensional Space}

\srcpage{284}

The present Appendix is devoted to the visual representation of those aspects of the
geometry of multi-dimensional spaces that are essential for the considerations set forth
in Chapter~VII when describing the methods of simultaneously transmitting on the picture
plane the ordinary and the mystical space.  The necessary clarity can be achieved only
by simplifying the problem under consideration.  We use here an approach that is very
commonly applied for this purpose in geometric reasoning.  As already noted in
Chapter~VII, a visual representation of three-dimensional ``layers'' situated in
four-dimensional space is impossible.  We therefore consider an analogous problem,
passing to simpler spaces---two-dimensional and three-dimensional.

Suppose there exist two-dimensional beings living in a two-dimensional space---for
instance, some ``flat'' creatures that move only on a certain plane and do not know
that there is a third dimension in the world, i.e.\ they cannot even imagine motion
perpendicular to the plane on which they live.  Obviously, the movements of these flat
beings are described by two coordinates.  In contrast to these creatures, we
(three-dimensional beings) can readily visualise that, in addition to the named plane,
there may be another plane, parallel to the first, on which flat beings can also live.
Denote the first plane~$\pi_1$ and the second~$\pi_2$.  Whatever the distance between
$\pi_1$ and~$\pi_2$, the beings inhabiting them cannot pass from one plane to the
other, since they are unable to perform the motion perpendicular to the planes that
this would require.

An important special case of this example is that in which the distance between~$\pi_1$
and~$\pi_2$ is infinitely small.  Most simply, one may imagine the planes~$\pi_1$
and~$\pi_2$ to coincide but to belong to different sides---say, $\pi_1$ corresponds to
the ``upper'' and $\pi_2$ to the ``lower'' surface of some infinitely thin plane.  Then
a situation arises in which two groups of heterogeneous beings live on one plane; these
groups never meet one another, they can, without interfering with each other, occupy the
same regions of ``their'' two-dimensional space, and so on---all because they are
displaced relative to one another by an infinitely small distance, but in a direction
that is incomprehensible to them---the ``third'' direction.

\srcpage{285}

The constructed example is in essence analogous to that considered in Chapter~VII, where
the discussion concerned two three-dimensional spaces infinitely slightly displaced
relative to one another in the ``fourth'' direction, allowing two ``lives''---ordinary
and mystical---to exist in the same volume (of a room, for instance).  Obviously the
analogies can be extended: mystical and ordinary space are not completely independent,
they are linked not materially but spiritually (by prayer, etc.), but this too can be
supposed of the flat beings inhabiting the sides~$\pi_1$ and~$\pi_2$ of the introduced
plane.  Further, if one allows that in certain exceptional cases the flat beings of~$\pi_1$
can ``seep through'' from the ``upper'' to the ``lower'' surface, they will possess the
ability to ``appear'' to the beings of~$\pi_2$ ``from nothing'', and so on.

For the problem under consideration here, what matters is not the analogies of the
properties of the complex of ordinary and mystical spaces and those of the two sides
($\pi_1$ and~$\pi_2$) of a single plane, but rather the rational methods of \emph{depicting}
such two spaces.  Begin with the vivid example of a two-sided plane.  Since a drawing
is always one-sided, the only way of simultaneously depicting the sides~$\pi_1$ and~$\pi_2$
on a drawing (without the depicted sides overlapping one another) is the \emph{method of
sections}: wherever the side~$\pi_1$ is depicted, the depiction of~$\pi_2$ must be
interrupted, and to facilitate their distinction the images of sides~$\pi_1$ and~$\pi_2$
should be painted in different colours.  Precisely the same procedure must be used when
depicting events occurring in two infinitely close three-dimensional ``layers'' of a
four-dimensional space: if one builds a visual three-dimensional model of such a space,
some regions of the model's volume must be assigned to conveying the ordinary space,
while in other regions of the model the mystical space is depicted.  As is clear from
what has been said, this is also the method of sections: two coexisting three-dimensional
spaces are shown alternately.

The transfer of such a three-dimensional model onto a picture plane (for instance, an
icon) no longer differs in any way from any other depiction of some three-dimensional
space, since the method of sections has reduced the task of conveying the geometry of two
three-dimensional ``layers'' of a four-dimensional space to the task of conveying the
geometry of a three-dimensional space in which sections of these ``layers'' neighbour one
another rather than ``overlapping'' one another.

\srcpage{286}

Of course, the method of sections is unable to show simultaneously both the ordinary and
the mystical spaces across the whole picture plane (just as one cannot show both sides of
a two-sided plane on one side of a drawing), but the possibility of simultaneously showing
both named spaces, even partially, is often more important than a complete depiction of
either one of these two spaces alone.
